\scanpage{56}
\mspage{V}

\uline{Cor} 1 Supposons que $\unsure{G|C}$ soit premier aux car. résiduelles. \deleted{Alors \unsure{la}} Supposons que pour tt sous-schéma affine, intègre normal $S_0$ de $S$, et tt \uline{\unsure{revêtmt|revêt|ouvert}} \uline{étale} \inserted{\unsure{intègre}} $S'$ de $S_0$, enfin tt \unsure{sous-gpe|sous-groupe} $H$ de $G\times G$, on ait $H^1(S',H)$ fini. Alors pour tt $X$ de t. f. sur $S$, $H^1(X,G)$ est fini.

(\unsure{Si|Or} $G$ est commutatif, il suffit de faire $H=G$ \deleted{\unsure{et}}, et on a une condition \unsure{néc.|née} et suffisante).

\uline{Corollaire} 2 Conditions équivalentes.

(i) Pour tt $X$ de t. f. sur $Y$, et tt \unsure{grpe|gpe|groupe} \inserted{fini} $G$ premier aux car. résiduelles de $X$, on a $H^1(X\unsure{,|.}G)=\text{fini}$

(ii) \deleted{\unsure{Pour}} \unsure{Idem|Item}, avec $X$ \unsure{revêtmt|revêt} \deleted{fini} étale intègre d'un sous-schéma affine normal intègre $S_0$ de $S$.

\uline{NB} Dans le cor. 1, on a utilisé le fait que si $k$ est un corps \unsure{alg.|alg} clos, alors pour tt $X$ de t. f. sur $k$ et $G$ premier à la car. de $k$, $H^1(X,G)$ est fini. Voyons comment ceci peut être prouvé par les méthodes envisagées ici.
% archive sheet number: 55
